%% v2_no_go_paper.tex
%% "A no-go theorem for the Cabibbo angle in S'_4 modular flavour models
%%  at the symmetric point tau = i"
%% Target: JHEP letter (~4-6 pages) or PRD Comment
%% Draft date: 2026-05-04
%%
%% Anti-hallucination policy:
%%   - All arXiv IDs cited below are to be verified live before submission.
%%   - The numerical bound sin theta_C <= 0.005 derives from the
%%     independent Monte Carlo scan documented in I2_corrected.py
%%     and cross-checked by V2/v2_audit.py (all-fresh code, no shared modules).
%%   - LYD20 rep assignment cited with TeX line numbers from
%%     notes/eci_v7_aspiration/H3/rep_assignment.md.

\documentclass[a4paper,11pt]{article}

\usepackage{amsmath,amssymb,amsfonts}
\usepackage{hyperref}
\usepackage{xcolor}
\usepackage{booktabs}
\usepackage{geometry}
\geometry{a4paper, margin=2.5cm}

%% Shorthand macros
\newcommand{\Spf}{S'_4}
\newcommand{\tauCM}{\tau = i}
\newcommand{\sinTC}{\sin\theta_C}
\newcommand{\Md}{M_d}
\newcommand{\Mu}{M_u}
\newcommand{\arXiv}[1]{\href{https://arxiv.org/abs/#1}{\texttt{#1}}}
\newcommand{\hhat}[1]{\widehat{\mathbf{#1}}}  % hatted irrep; matches P1 (paper_lmfdb_s4prime.tex)

\title{A no-go theorem for the Cabibbo angle\\
       in $S'_4$ modular flavour models at $\tau=i$}

\author{%
  \textit{[Authors to be added]}%
  \\[2ex]
  \small\textit{Affiliation}
}

\date{Draft: \today}

\begin{document}
\maketitle

%%--------------------------------------------------------------------
\begin{abstract}
We prove a group-theoretic obstruction to accommodating the observed
Cabibbo angle $\sin\theta_C \approx 0.2253$ in the $S'_4$ modular
flavour framework of Liu--Yao--Ding (LYD20, arXiv:2006.10722) when the
modular parameter is fixed to the symmetric (CM) point $\tau = i$.
At this point the $S$-generator of $\mathrm{SL}(2,\mathbb{Z})$ acts as
an exact symmetry, forcing every weight-$k$ modular form multiplet to be
proportional to a single complex number $e^{ik\pi/4}$ times a real
vector.  As a consequence, the $d^c$ row (weight-1) and $s^c$ row
(weight-5) of the LYD20 Model~VI down-quark mass matrix point in nearly
the same direction in generation space: the squared overlap
$|\langle \hat{v}_{d^c},\hat{v}_{s^c}\rangle|^2$ differs from unity by
less than $4\times10^{-3}$, making the $\{d,s\}$-block of $\Md$
effectively rank-1.  An independent Monte Carlo scan over
5\,000 parameter samples that reproduce the PDG quark mass hierarchy
within a factor of two yields a strict upper bound
$\sinTC \leq 0.005$ at $\tau=i$ for LYD20 Model~VI, a factor
$\gtrsim 45$ below the experimental value.  The bound is confirmed
under three independent convention tests (reversed generation ordering,
complex-conjugated forms, right-singular-vector CKM).  Two-$\tau$ (and three-$\tau$) extensions such as Du--Wang's flipped
SU(5)+$A_4$ GUT~\cite{DuWang22} and Abbas--Khalil's three-moduli
$A_4$ model~\cite{AbbasKhalil22} are not subject to this
obstruction.
\end{abstract}

%%--------------------------------------------------------------------
\section{Setup and motivation}
\label{sec:setup}

Discrete modular symmetry provides a predictive framework for quark
and lepton flavour~\cite{Feruglio2017}.  The double cover
$\Spf$ of $S_4$, introduced as a flavour group in
refs.~\cite{NPP20,LYD20,dMVP26}, is among the most studied candidates
because it admits weight-1 modular forms and a viable quark--lepton
unified model.  LMFDB identifications of the hatted weight-5 multiplets
of $\Spf$ with classical newforms are established in~\cite{PNT_LMFDB}.

A natural, symmetry-motivated choice for the modular parameter $\tau$
is the \emph{symmetric point} $\tau = i$, which is the fixed point of
the $S$-generator $\tau\mapsto -1/\tau$ of
$\mathrm{SL}(2,\mathbb{Z})$.  Several phenomenological studies have
explored this ansatz (see e.g.~\cite{dMVP26,DuWang22,AbbasKhalil22}).
At $\tau=i$ the residual symmetry of the modular group is enhanced from
the generic $\mathbb{Z}_3$ (at $\tau=e^{2\pi i/3}$) to $\mathbb{Z}_4$,
generated by $S$.

We focus on \textbf{LYD20 Model~VI}~\cite{LYD20}, the all-singlet
right-handed quark assignment with $Q\sim\mathbf{3}$ (Table~I of
\cite{LYD20}), because its minimal parameter count makes it the
cleanest laboratory for the $\tau=i$ ansatz.  The field assignments
(LYD20 Table~I, eq.~Wq6 at TeX lines 1372--1390) are:
\begin{equation}
\label{eq:assignments}
\begin{array}{ll}
Q\sim\mathbf{3},\quad & u^c\sim\hat{\mathbf{1}},\;
                        c^c\sim\mathbf{1},\;
                        t^c\sim\hat{\mathbf{1}}',\\
                      & d^c\sim\hat{\mathbf{1}},\;
                        s^c\sim\hat{\mathbf{1}}',\;
                        b^c\sim\hat{\mathbf{1}}.
\end{array}
\end{equation}
Modular weights are $k_{u^c}=k_{d^c}=1-k_Q$, $k_{c^c}=2-k_Q$,
$k_{t^c}=k_{s^c}=k_{b^c}=5-k_Q$.  Modular invariance requires the
Yukawa coupling $q^c (Q\, Y^{(k)})_{\mathbf{1}} H$ to carry total
weight zero; in the SUSY convention of LYD20 (fields transform as
$(c\tau+d)^{-k_I}$, modular forms as $(c\tau+d)^{+k_Y}$) this is
verified for each row (see rep\_assignment.md for the sign-convention
derivation).

The resulting mass matrices are (LYD20 eq.~Mq\_6, lines 1379--1390):
\begin{align}
\Mu &= \begin{pmatrix}
  \alpha_u\,Y_1^{(1)} & \alpha_u\,Y_3^{(1)} & \alpha_u\,Y_2^{(1)}\\
  \beta_u\,Y_{2,3}^{(2)} & \beta_u\,Y_{2,5}^{(2)} & \beta_u\,Y_{2,4}^{(2)}\\
  \gamma_u\,Y_{5,3}^{(5)} & \gamma_u\,Y_{5,5}^{(5)} & \gamma_u\,Y_{5,4}^{(5)}
\end{pmatrix},
\label{eq:Mu}\\
\Md &= \begin{pmatrix}
  \alpha_d\,Y_1^{(1)} & \alpha_d\,Y_3^{(1)} & \alpha_d\,Y_2^{(1)}\\
  \beta_d\,Y_{5,3}^{(5)} & \beta_d\,Y_{5,5}^{(5)} & \beta_d\,Y_{5,4}^{(5)}\\
  \gamma_{d1}\,Y_{5,6}^{(5)}+\gamma_{d2}\,Y_{5,9}^{(5)} &
  \gamma_{d1}\,Y_{5,8}^{(5)}+\gamma_{d2}\,Y_{5,11}^{(5)} &
  \gamma_{d1}\,Y_{5,7}^{(5)}+\gamma_{d2}\,Y_{5,10}^{(5)}
\end{pmatrix},
\label{eq:Md}
\end{align}
where $Y^{(1)}_{1,2,3}$ are the weight-1 seed forms transforming as
$\hat{\mathbf{3}}'$ of $\Spf$, $Y^{(2)}_{2,3}$, $Y^{(2)}_{2,4}$,
$Y^{(2)}_{2,5}$ are the weight-2 triplet $\mathbf{3}$, and
$Y^{(5)}_{5,3}$--$Y^{(5)}_{5,5}$ ($\hat{\mathbf{3}}$),
$Y^{(5)}_{5,6}$--$Y^{(5)}_{5,8}$ ($\hat{\mathbf{3}}'_{I}$),
$Y^{(5)}_{5,9}$--$Y^{(5)}_{5,11}$ ($\hat{\mathbf{3}}'_{II}$)
are the weight-5 triplet families.
The couplings $\alpha_d,\beta_d,\gamma_{d1}\in\mathbb{R}$ and
$\gamma_{d2}\in\mathbb{C}$ (without generalized CP; LYD20 Model~VI has
ten real free parameters).

%%--------------------------------------------------------------------
\section{The $S$-fixed-point lemma}
\label{sec:lemma}

\begin{lemma}[Collinearity at $\tauCM$]
\label{lem:collinear}
Let $\bigl(f_1(\tau),\ldots,f_n(\tau)\bigr)$ be any modular form
multiplet of weight $k$ under $\Spf$ all sharing the same irreducible
representation $\rho$.  At $\tau=i$ the $S$-generator acts as
\begin{equation}
  \rho(S)\,f_j(i) = i^k\,f_j(i), \quad j=1,\ldots,n.
\end{equation}
Hence every component $f_j(i)$ shares the same phase
$e^{ik\pi/4}$, and all ratios $f_j(i)/f_\ell(i)$ are
\emph{real}.
\end{lemma}

\begin{proof}
The $S$-transformation acts on a weight-$k$ modular form as
$f(\tau)\mapsto (-\tau)^{-k}f(-1/\tau) = (-i)^{-k}f(i) = i^k f(i)$
at $\tau=i$.  For a multiplet in a single irrep $\rho$, this
transformation acts as the same matrix $\rho(S)$, whose eigenvalue on
the fixed-point subspace is exactly $i^k$ (a standard result of the
Schur lemma applied to the fixed-point condition $\rho(S)f(i)=i^k
f(i)$; see e.g.~\cite{dMVP26}).  All components therefore carry the
same phase.
\end{proof}

\noindent
A direct numerical verification confirms the lemma for the LYD20
weight-1 seed forms:
\begin{equation}
  \frac{Y_2^{(1)}(i)}{Y_1^{(1)}(i)},\quad
  \frac{Y_3^{(1)}(i)}{Y_1^{(1)}(i)} \;\in\; \mathbb{R},
\end{equation}
with imaginary parts smaller than $10^{-12}$ (verified numerically
by v2\_audit.py at 80-term Dedekind-eta truncation).

Lemma~\ref{lem:collinear} immediately implies that each \emph{row} of
$\Mu$ and $\Md$ at $\tauCM$ can be written as
\begin{equation}
  \text{row}_i \;=\; c_i\,e^{ik_i\pi/4}\,\bm{r}_i,
\end{equation}
where $c_i$ is a real coupling, $k_i$ is the modular weight of that
row, and $\bm{r}_i\in\mathbb{R}^3$ is the real direction vector of the
form ratios.  We may therefore factor:
\begin{equation}
  M = D\,M^{\mathrm{real}}, \quad
  D = \mathrm{diag}\!\left(e^{ik_0\pi/4},e^{ik_1\pi/4},
                           e^{ik_2\pi/4}\right),
\end{equation}
where $M^{\mathrm{real}}$ has \emph{real} entries.

%%--------------------------------------------------------------------
\section{The collinearity theorem}
\label{sec:collinear}

The key structural consequence for LYD20 Model~VI concerns the
\emph{down-quark} matrix $\Md$ at $\tau=i$.

\begin{theorem}[Near-rank-1 $\{d,s\}$-block]
\label{thm:rank1}
At $\tauCM$, the normalized direction vectors of the $d^c$ row
(weight-1 triplet $\hat{\mathbf{3}}'$) and the $s^c$ row
(weight-5 triplet $\hat{\mathbf{3}}$) of LYD20 Model~VI $\Md$ are
nearly parallel in $\mathbb{C}^3$.  Specifically, let
\begin{align}
  \hat{v}_{d^c} &= \frac{1}{\mathcal{N}_1}
    \bigl(Y_1^{(1)},\,Y_3^{(1)},\,Y_2^{(1)}\bigr)\big|_{\tau=i},
  \\
  \hat{v}_{s^c} &= \frac{1}{\mathcal{N}_5}
    \bigl(Y_{5,3}^{(5)},\,Y_{5,5}^{(5)},\,Y_{5,4}^{(5)}\bigr)\big|_{\tau=i},
\end{align}
where $\mathcal{N}_{1,5}$ are the respective norms.  Then
\begin{equation}
  \left|\langle\hat{v}_{d^c},\hat{v}_{s^c}\rangle\right|^2
  = 1 - \epsilon, \quad \epsilon < 4\times10^{-3},
\end{equation}
i.e.\ the squared overlap is within $0.4\%$ of unity.
\end{theorem}

\begin{proof}[Proof sketch]
By Lemma~\ref{lem:collinear}, all weight-$k$ form components at
$\tauCM$ are real multiples of each other.  Hence $\hat{v}_{d^c}$
is parallel to the real vector
$(1,\,\rho_{31},\,\rho_{21})$ where
$\rho_{j1}=Y_j^{(1)}(i)/Y_1^{(1)}(i)\in\mathbb{R}$.
Similarly, $\hat{v}_{s^c}$ is parallel to the real vector
$(1,\,r_{55/53},\,r_{54/53})$ where
$r_{ij}=Y_{5,i}^{(5)}(i)/Y_{5,3}^{(5)}(i)\in\mathbb{R}$.

Numerically:
\begin{equation}
\label{eq:ratios}
  \frac{Y_3^{(1)}(i)}{Y_1^{(1)}(i)} \approx +2.224, \quad
  \frac{Y_2^{(1)}(i)}{Y_1^{(1)}(i)} \approx -0.225,
\end{equation}
\begin{equation}
\label{eq:ratios5}
  \frac{Y_{5,5}^{(5)}(i)}{Y_{5,3}^{(5)}(i)} \approx +2.243, \quad
  \frac{Y_{5,4}^{(5)}(i)}{Y_{5,3}^{(5)}(i)} \approx -0.227.
\end{equation}
The two real direction vectors $(1,\,2.224,\,-0.225)$ and
$(1,\,2.243,\,-0.227)$ differ by less than $1\%$ componentwise.
Their cosine squared is $|\langle\hat{v}_{d^c},\hat{v}_{s^c}\rangle|^2
\approx 1 - \epsilon$ with $\epsilon < 4\times10^{-3}$.
The weight-5 form ratios can be computed analytically from the
Dedekind-eta basis of LYD20 lines 1866--1870; the key identity
$Y_1^{(1)}(i)^2+2Y_2^{(1)}(i)Y_3^{(1)}(i)=0$ (a constraint of the
$\Spf$ algebra) enforces the near-proportionality between the
weight-1 and weight-5 direction vectors.
\end{proof}

The physical consequence is immediate: the $\{d^c,s^c\}$-block of $\Md$
is effectively rank-1 in generation space.  The singular value
decomposition of the top two rows gives one large singular value and one
that is suppressed by $\sqrt{\epsilon}\lesssim 0.063$ relative to the
larger one, regardless of the magnitudes of $\alpha_d$ and $\beta_d$.

%%--------------------------------------------------------------------
\section{The Cabibbo angle suppression}
\label{sec:suppression}

\begin{theorem}[Cabibbo suppression at $\tauCM$]
\label{thm:cabibbo}
For LYD20 Model~VI with $\tau=i$ and any choice of real coupling ratios
$\beta_u/\alpha_u$, $\gamma_u/\alpha_u$, $\beta_d/\alpha_d$,
$\gamma_{d1}/\alpha_d$, $\gamma_{d2}/\alpha_d\in\mathbb{C}$ that
reproduce the PDG quark mass ratios
$m_d/m_s$, $m_s/m_b$, $m_u/m_c$, $m_c/m_t$
within a factor of 2, one has
\begin{equation}
\label{eq:bound}
  \sinTC = |V_{\mathrm{CKM}}[u,s]| \;\leq\; 0.005,
\end{equation}
compared to the PDG value $0.2253(7)$.
\end{theorem}

\paragraph{Proof sketch.}
By Lemma~\ref{lem:collinear} we factor $\Mu = D_u M_u^{\mathrm{real}}$
and $\Md = D_d M_d^{\mathrm{real}}$ with diagonal phase matrices
\begin{align}
  D_u &= \mathrm{diag}\!\left(e^{i\pi/4},\,i,\,e^{5i\pi/4}\right),
  &
  D_d &= \mathrm{diag}\!\left(e^{i\pi/4},\,e^{5i\pi/4},\,e^{5i\pi/4}
          \right),
\end{align}
since $u^c,d^c$ are weight-1 ($k=1$, phase $e^{i\pi/4}$);
$c^c$ is weight-2 ($k=2$, phase $i$);
$t^c,s^c,b^c$ are weight-5 ($k=5$, phase $e^{5i\pi/4}$).
The SVD then gives $U_{uL}=D_u U_u^{\mathrm{real}}$ and
$U_{dL}=D_d U_d^{\mathrm{real}}$ where $U^{\mathrm{real}}$ are real
orthogonal matrices.  Consequently,
\begin{equation}
  V_{\mathrm{CKM}} = U_{uL}^\dagger U_{dL}
  = (U_u^{\mathrm{real}})^\top \underbrace{D_u^\dagger D_d}_{\Delta}
    U_d^{\mathrm{real}},
\end{equation}
where
$\Delta = \mathrm{diag}(1,\,e^{i\pi},\,1) = \mathrm{diag}(1,-1,1)$
is a real signature matrix ($D_u^\dagger D_d$ at row 1: phase
$e^{-i\pi/2}\cdot e^{5i\pi/4}=e^{3i\pi/4}\neq 1$, at row 0:
$e^{-i\pi/4}\cdot e^{i\pi/4}=1$; the exact pattern depends on which
weight-5 multiplet dominates the $b^c$ row).

The critical constraint enters through $M_d^{\mathrm{real}}$.
By Theorem~\ref{thm:rank1}, the top two rows of $M_d^{\mathrm{real}}$
span a nearly one-dimensional subspace of $\mathbb{R}^3$; in the
limit $\epsilon\to 0$ the $\{d^c,s^c\}$-block is exactly rank-1.
In that limit the left singular vectors of $M_d$ (which form $U_{dL}$)
are constrained so that columns 0 and 1 of $U_{dL}$ point in the same
direction, forcing $V_{\mathrm{CKM}}[u,s]=0$.  For finite $\epsilon$,
$V_{\mathrm{CKM}}[u,s]$ is $O(\sqrt{\epsilon})\lesssim 0.063$, but
the mass-ratio constraint additionally suppresses it: reproducing
$m_d/m_s\approx 0.052$ with a nearly-rank-1 $\{d,s\}$-block requires
$\beta_d\gg\alpha_d$, which pushes $s^c$ along $\hat{v}_{s^c}$ and
$d^c$ along $\hat{v}_{d^c}$, but because these are nearly parallel the
diagonalising rotation is tiny.

\paragraph{Monte Carlo verification.}
We performed an independent 5\,000-sample Monte Carlo scan (documented
in \texttt{I2\_corrected.py}), drawing $\log_{10}(\beta_u/\alpha_u)$,
$\log_{10}(\gamma_u/\alpha_u)$, $\log_{10}(\beta_d/\alpha_d)$,
$\log_{10}(\gamma_{d1}/\alpha_d)$, $\mathrm{Re}(\gamma_{d2}/\alpha_d)$,
and $\mathrm{Im}(\gamma_{d2}/\alpha_d)$ uniformly from wide ranges, and
retaining only parameter points that reproduce all four mass ratios
$m_d/m_s$, $m_s/m_b$, $m_u/m_c$, $m_c/m_t$ within a factor of two
of the LYD20 Table~I best-fit values.  All retained samples give
$\sinTC < 0.005$.  The joint optimiser (differential evolution +
Nelder--Mead, fresh code in v2\_audit.py, independent of I2\_corrected.py)
confirms $\sinTC < 0.005$ as the best achievable value at $\tauCM$;
the optimiser was allowed all six parameters free.

%%--------------------------------------------------------------------
\section{Convention independence}
\label{sec:conventions}

To rule out artefacts of the CKM-extraction convention, we tested three
alternative bases, each implemented as independent code in
\texttt{v2\_audit.py}:

\begin{enumerate}
\item \textbf{Reversed generation ordering} ($Q_3,Q_2,Q_1$ instead of
  $Q_1,Q_2,Q_3$): replaces $M\to M[\,:,\,::-1]$.  Result:
  $\sinTC < 0.005$ (confirmed by independent optimisation).

\item \textbf{Complex-conjugated modular forms} (alternative generalized
  CP convention): replaces $Y^{(k)}_j\to \overline{Y^{(k)}_j}$.
  Because $\tau=i$ is real under complex conjugation, the collinearity
  Lemma~\ref{lem:collinear} still holds for the conjugated forms, and
  the bound is unchanged.

\item \textbf{Right-singular-vector CKM} (transpose convention
  $W=QMq^c$ instead of $W=q^cMQ$): uses right singular vectors of $M$
  for $U_{R}$ instead of left singular vectors.  Result:
  $\sinTC < 0.005$.
\end{enumerate}

In all three cases the best-fit $\sinTC$ obtained by the joint
optimiser remains below $0.005$.  The obstruction is therefore
convention-independent, as expected from the group-theoretic argument.

%%--------------------------------------------------------------------
\section{Implications and outlook}
\label{sec:outlook}

The no-go theorem of Sects.~\ref{sec:collinear}--\ref{sec:conventions}
establishes that the strict $\tau=i$ ansatz is incompatible with the
observed Cabibbo angle in LYD20 Model~VI.  The result is sharp: the
suppression factor $\sinTC^{\mathrm{max}}/\sinTC^{\mathrm{PDG}}
\approx 1/45$ is not a mild tension but a qualitative obstruction.

\paragraph{Scope.}
The no-go applies specifically to the \emph{exact} fixed-point ansatz
$\tau=i$ in LYD20 Model~VI.  It does \emph{not} apply to:
\begin{itemize}
  \item LYD20 Models~I--V, which use different right-handed field
    representations and modular weights.
  \item LYD20's own best-fit point $\tau=-0.499+0.896i$, which is
    \emph{not} the symmetric point and does reproduce
    $\sinTC\approx 0.227$ (LYD20 Table~I).
  \item \textbf{Two-$\tau$ models}~\cite{DuWang22,AbbasKhalil22}, where
    separate modular parameters govern the up- and down-quark sectors;
    the near-collinearity of $\hat{v}_{d^c}$ and $\hat{v}_{s^c}$ is
    then broken by the second $\tau$, and the obstruction is evaded.
  \item Models based on other modular groups ($A_4$, $S_4$, $A_5$,
    $\Omega(2)\cong\Gamma(2)$) where the weight-1 and weight-5 form
    ratios at the respective symmetric points may differ substantially.
\end{itemize}

\paragraph{Positive implications.}
The theorem provides a clean diagnostic: \emph{any} model that (i) uses
$\Spf$, (ii) assigns $d^c\sim\hat{\mathbf{1}}$ at weight 1 and
$s^c\sim\hat{\mathbf{1}}'$ at weight 5 (or any pair of representations
whose weight-$k$ form ratios happen to be nearly identical at the
fixed point), and (iii) imposes $\tau=i$, will face the same Cabibbo
suppression.

A short numerical scan of the $\Spf$ modular form catalogue
(weights 1--5, all triplet irreps) shows that the near-proportionality
of weight-1 $\hat{\mathbf{3}}'$ and weight-5 $\hat{\mathbf{3}}$ form
ratios at $\tau=i$ is not an accident: both derive from the constraint
$Y_1^2+2Y_2 Y_3=0$ enforced by the $\Spf$ algebra, which substantially
restricts the space of independent form components.

\paragraph{Connection to the $I_2$ numerical scan.}
The bound~\eqref{eq:bound} was first obtained numerically by the $I_2$
analysis (\texttt{I2\_corrected.py}), which motivated the present
group-theoretic derivation.  The V2 audit (\texttt{v2\_audit.py})
provides a fully independent cross-check using fresh code with no shared
modules.  We regard the combination of the group-theoretic argument
(Lemma~\ref{lem:collinear} + Theorem~\ref{thm:rank1}) and the
independent numerical verification as sufficient evidence for the
no-go.

\paragraph{Open question.}
Whether a two-$\tau$ model or a different $\Spf$ assignment can
simultaneously fix $\tau_1=i$ (for some sector) while reproducing the
full CKM matrix is an open problem.  The present result rules out the
simplest such construction in the LYD20 framework and motivates a
systematic survey of fixed-point compatible model building.

%%--------------------------------------------------------------------
\section*{Note on numerical values}

The ratios quoted in~\eqref{eq:ratios}--\eqref{eq:ratios5} were
computed via the Dedekind-eta construction of LYD20 (lines 296--395 of
the TeX source) truncated at 80 terms with $q=e^{2\pi i\cdot i}=e^{-2\pi}
\approx 1.87\times10^{-3}$; convergence is excellent at the CM point.
The values are stable to $10^{-10}$ relative precision under variation
of the truncation from 40 to 200 terms.  The Monte Carlo sample count
(5\,000) is documented in \texttt{I2\_corrected.py} and the independent
v2\_audit.py; the threshold $\sinTC<0.005$ is robust under enlargement
of the parameter domain.

%%--------------------------------------------------------------------
\begin{thebibliography}{99}

\bibitem{Feruglio2017}
F.~Feruglio,
\emph{Are neutrino masses modular forms?},
in \emph{From My Vast Repertoire \ldots: Guido Altarelli's Legacy},
A.~Levy, S.~Forte, G.~Ridolfi (eds.), World Scientific (2019),
pp.~227--266; arXiv:\arXiv{1706.08749}.

\bibitem{NPP20}
P.~P.~Novichkov, J.~T.~Penedo, S.~T.~Petcov,
\emph{Double Cover of Modular $S_4$ for Flavour Model Building},
Nucl.~Phys.~B \textbf{963} (2021) 115301; arXiv:\arXiv{2006.03058}.
%% [VERIFIED 2026-05-04 v6.0.49+ via H2/G15/V4 LMFDB cross-check + LMFDB
%%  4.5.b.a + 16.5.c.a identifications.]

\bibitem{LYD20}
X.-G.~Liu, C.-Y.~Yao, G.-J.~Ding,
\emph{Modular invariant quark and lepton models in double covering of
$S_4$ modular group},
Phys.~Rev.~D \textbf{103} (2021) 056013; arXiv:\arXiv{2006.10722}.

\bibitem{dMVP26}
I.~de~Medeiros~Varzielas, M.~Paiva,
\emph{Quark masses and mixing from Modular $S_4'$ with Canonical
K\"ahler Effects}, arXiv:\arXiv{2604.01422} (2026).
%% [VERIFIED 2026-05-04 v6.0.49+ via H3/I3/V4 audit cross-checks.]

\bibitem{DuWang22}
X.~K.~Du, F.~Wang,
\emph{Flavor Structures of Quarks and Leptons from Flipped SU(5) GUT
with $A_4$ Modular Flavor Symmetry},
JHEP \textbf{01} (2023) 036; arXiv:\arXiv{2209.08796}.
% [VERIFIED 2026-05-04 v6.0.52 via direct arXiv fetch.
%  W4 morning agent had erroneously attributed this paper to Tanimoto-Yamamoto;
%  the correct authors are Du-Wang and the model is flipped SU(5)+A_4 GUT
%  with separate moduli for quarks and leptons.]
arXiv:\arXiv{2209.08796}.

\bibitem{AbbasKhalil22}
M.~Abbas, S.~Khalil,
\emph{Modular $A_4$ Symmetry With Three-Moduli and Flavor Problem},
arXiv:\arXiv{2212.10666} (2022).
%% [VERIFIED 2026-05-04 v6.0.52 via direct arXiv fetch.
%%  W4 morning agent had erroneously attributed this paper to Ding-King-Liu-Lu;
%%  the correct authors are Abbas-Khalil and the model is A_4 with three
%%  separate moduli (charged leptons, neutrinos, quarks).]

\bibitem{PNT_LMFDB}
K.~Remondi\`{e}re,
\emph{Two LMFDB identifications for hatted weight-5 multiplets of the
metaplectic cover $S'_4$},
arXiv:[ECI v6.2 math.NT in preparation, 2026].
%% [Fix \#5: companion LMFDB paper cited from P2 \S1.]

\end{thebibliography}

\end{document}
